\section{Introduction}

Leaderboards for agent benchmarks are read as measurements. The gap between
adjacent entries is cited as evidence that one system is better than another,
and differences of two or three points routinely carry that weight. On the
board we study, 37 of 40 entries are the result of a single run, and no entry
reports an interval.

A benchmark that drives a real operating system through a real user interface
against real network services will have run-to-run variance. That is the price
of not being a static test set, and it is a price worth paying: the alternative
measures something much less like deployment. What has been missing is a
number for it, and a way to tell when the number has changed.

We pinned every control the tooling of
MobileWorld~\citep{kong2026mobileworld} exposes and ran ten repetitions over 29
GUI-only tasks, 290 runs in total. We set out to measure a noise floor. We
found something larger: the configuration was not, in the event, fixed, and
nothing in the tooling or the reporting format would have told us.

\paragraph{The claim.} A benchmark is an instrument for measuring agent
capability. This instrument moves while nobody is recording, and current
reporting practice contains no mechanism that would reveal it.

\paragraph{Contributions.}
\begin{itemize}\itemsep2pt
\item \textbf{A methodological error in determinism checks.} Identical requests
  at \texttt{temperature{=}0} with a fixed seed returned different completions
  on both endpoints we tested, though both advertise seed support. An earlier
  check of ours with a short prompt returned byte-identical output twice and
  looked like proof of the opposite. Under greedy decoding a thirty-token answer
  comes back byte-identical almost regardless of the endpoint; determinism has
  to be probed at the output length the real workload produces.
\item \textbf{A documented mid-run serving change, with a detector.} Between
  repetition 3 and 4, with every flag unchanged, ten tasks got worse and none
  got better ($p = 0.002$), and median completion tokens per action fell 20.1\%
  in one step. We give five independent lines of evidence and the cheap signal
  that exposes it.
\item \textbf{A noise floor for this configuration, and what it implies.}
  Within one serving configuration, two single runs need to differ by
  8.15 points before the gap exceeds what repetition alone produces. The figure
  is a property of a resampling model fitted to one agent, one endpoint and one
  reweighted task sample, not a constant of the benchmark.
\item \textbf{Findings that need no compute.} Eight properties of the public
  repository, each with where to look, that bear on whether two entries
  are comparable at all. A reader can check each one without trusting any
  measurement of ours.
\end{itemize}

A benchmark that reports only a scalar score is a measuring device with no
self-check: it cannot distinguish a system that improved from an instrument
that moved, and neither can the metadata recorded beside it. Treating the
harness as a verification system --- recording the signals a failure would show
up in, not only the number --- is what made Section~\ref{sec:endpoint} visible
at all, at a cost of one integer per request.

The first three compose into one argument rather than three results, and the
fourth asks for no trust at all. Determinism controls are why one must pin the
endpoint; the endpoint change is why pinning is not enough and why nothing would
have told us; the noise floor prices the consequence.
